\ifx\allfiles\undefined
\documentclass[12pt, a4paper, oneside, UTF8]{ctexbook}
\def\path{../config}
\input{\path/package.tex}

% 定理环境设置
\input{\path/theorem1_zh.tex}

\input{\path/custom.tex}

% 修改参数改变封面样式，0 默认原始封面、内置其他1、2、3种封面样式
\def\myIndex{3}


\ifnum\myIndex>0
    \input{\path/cover_package_\myIndex} 
\fi

\def\myTitle{高等数学笔记}
\def\myAuthor{M.K.}
\def\myDateCover{\today}
\def\myDateForeword{\today}
\def\myForeword{前言}
\def\myForewordText{
    
    这是一个基于\LaTeX{}模板的高数笔记．
    内容参考了包括但不限于以下资料：
    \begin{itemize}
        \item 《高等数学》（同济大学出版社）第七版
        \item 《吉米多维奇高等数学习题精选精讲》（山东科学技术出版社）第二版
        \item 《高等数学作业集》（哈尔滨工业大学出版社）第一版
        \item \href{https://space.bilibili.com/1035929235}{B站UP主：一高数}的高数课程
        \item \href{https://space.bilibili.com/42428180}{B站UP主：考研竞赛数学}的试题讲解与高数课程
        \item \href{https://space.bilibili.com/391930545}{B站UP主：Maki的完美算数教室}的高数内容

    \end{itemize}

    封面来源于\href{https://www.pixiv.net/users/3952}{おにねこ}的作品\href{https://www.pixiv.net/artworks/136551599}{雨雫とプレアデス}
    如有侵权请联系我删除．
}
\def\mySubheading{副标题}


\begin{document}
% \input{../config/cover}
\else
\fi
\chapter{偏微分方程初步与全微分方程}

以自然界中的物理现象为引，本章我们将探讨多变量微积分的核心应用：全微分方程及偏微分方程初步。经典偏微分教材常常将这类理论放置在流体力学、热力学或电磁学的情境下进行直观解释。

\section{全微分方程~(Exact Differential Equations)}
\subsection{全微分的物理含义与 $\dfrac{\partial P}{\partial y}=\dfrac{\partial Q}{\partial x}$ 判定}

\begin{definition}[全微分方程及其含义]
    设一阶微分方程的形式为 $P(x,y)\d x + Q(x,y)\d y = 0$．如果在某区域内存在一个势函数 $u(x,y)$，使得其全微分恰好等于方程的左端：
    \[ \d u = \frac{\partial u}{\partial x}\d x + \frac{\partial u}{\partial y}\d y = P(x,y)\d x + Q(x,y)\d y \]
    则称该方程为\textbf{全微分方程}。
    
    \textbf{物理含义}上，这等价于平面向量场 $\vec{F} = (P, Q)$ 是一个保守场（无散/无旋场），且无耗散作功。方程的通解直接由等势线给出：$u(x,y) = C$．
\end{definition}

\begin{theorem}
    设函数 $P(x,y)$ 与 $Q(x,y)$ 在单连通区域内具有连续的一阶偏导数，则 $P\d x + Q\d y = 0$ 是全微分方程的充分必要条件为（即等价于标量旋度为零）：
    \[ \frac{\partial P}{\partial y} = \frac{\partial Q}{\partial x} \]
\end{theorem}

\subsection{偏积分解法~(Method of Partial Integration)}
如果方程满足上述判定条件 $\frac{\partial P}{\partial y} = \frac{\partial Q}{\partial x}$，那么必定存在势函数 $u(x,y)$ 满足 $\frac{\partial u}{\partial x} = P$ 且 $\frac{\partial u}{\partial y} = Q$。

要系统地复原出 $u(x,y)$，最严谨的代数手段是\textbf{偏积分}，其核心逻辑包含三步：
\begin{enumerate}
    \item \textbf{选向偏积分：} 利用 $\frac{\partial u}{\partial x} = P(x,y)$，暂时将 $y$ 视作常数，对 $x$ 进行积分：
    \[ u(x,y) = \int P(x,y) \d x + \varphi(y) \]
    【核心注意】：由于这是对 $x$ 的偏向不定积分，所以产生的“积分常数”实际上是一个可能包含另一变量 $y$ 的未知函数 $\varphi(y)$。
    \item \textbf{求导做比对：} 利用尚未使用的条件 $\frac{\partial u}{\partial y} = Q(x,y)$，将上式对 $y$ 求偏导，并强制等于 $Q(x,y)$：
    \[ \frac{\partial}{\partial y} \left( \int P(x,y) \d x \right) + \varphi'(y) = Q(x,y) \]
    \item \textbf{得出原函数：} 移项去解 $\varphi'(y)$。依据判定条件，此时等式中残留的 $x$ 会奇迹般地全部抵消掉！算出纯关于 $y$ 的 $\varphi'(y)$ 后，对 $y$ 积分求得 $\varphi(y)$。代回第一步即可拼凑出通解 $u(x,y) = C$。
\end{enumerate}
这就是求解全微分方程最基础、普适的“系统级保底解法”，这种方法也在物理中用于由保守力场反推标量势。

\subsection{常用全微分公式与配凑法}
除了老实地求解偏积分或者沿着路径进行曲线积分外，物理直觉上最常用的进阶解法是\textbf{配凑法 (Method of Grouping / Inspection)}。这要求我们非常熟悉以下常见的全微分形式（即微分乘除法则的逆推）：

\begin{table}[H]
    \centering
    \caption{常用全微分公式}
    \renewcommand\arraystretch{1.8}
    \begin{tabular}{ll}
        \hline
        \textbf{拼凑结构式 (Grouping)} & \textbf{原函数的全微分 (Exact Form)} \\
        \hline
        $y \d x + x \d y$ & $\d(xy)$ \\
        $\dfrac{x \d y - y \d x}{x^2}$ & $\d\left(\dfrac{y}{x}\right)$ \\
        $\dfrac{y \d x - x \d y}{y^2}$ & $\d\left(\dfrac{x}{y}\right)$ \\
        $\dfrac{x \d x + y \d y}{\sqrt{x^2+y^2}}$ & $\d(\sqrt{x^2+y^2})$ \\
        $\dfrac{x \d x + y \d y}{x^2+y^2}$ & $\d\left(\dfrac{1}{2}\ln(x^2+y^2)\right)$ \\
        $\dfrac{x \d y - y \d x}{x^2+y^2}$ & $\d\left(\arctan\dfrac{y}{x}\right)$ \\
        \hline
    \end{tabular}
\end{table}

\begin{example}
    求解一阶微分方程：$(3x^2+2y)\d x + (2x-3y^2)\d y = 0$．
\end{example}
\begin{solution}
    \textbf{方法一（检验法与偏积分解法）：}\\
    分析全微分方程结构，记 $P(x,y) = 3x^2+2y,\, Q(x,y) = 2x-3y^2$。\\
    由于 $P(x,y)$ 与 $Q(x,y)$ 及其偏导数在整个 $xy$-平面内均连续：\\
    $\dfrac{\partial P}{\partial y} = \dfrac{\partial}{\partial y}(3x^2+2y) = 2$\\
    $\dfrac{\partial Q}{\partial x} = \dfrac{\partial}{\partial x}(2x-3y^2) = 2$\\
    显然 $\dfrac{\partial P}{\partial y} \equiv \dfrac{\partial Q}{\partial x}$，所以原方程确实是一个确切的（Exact）全微分方程，必定存在标量势函数 $u(x,y)$。

    现在通过偏积分寻找 $u(x,y)$：
    由全微分定义，$\dfrac{\partial u}{\partial x} = P(x,y) = 3x^2+2y$。\\
    对 $x$ 进行第一轮偏偏积分（将 $y$ 视作未变化的参数），可以得到含有未知单变量函数 $\varphi(y)$ 的初始形式：
    \[ u(x,y) = \int (3x^2+2y) \d x + \varphi(y) = x^3 + 2xy + \varphi(y) \]
    接着，使用该结果去匹配本应满足的另外一个偏导条件 $\dfrac{\partial u}{\partial y} = Q(x,y)$：
    将刚刚得到的 $u(x,y)$ 对 $y$ 求偏导：
    \[ \frac{\partial u}{\partial y} = 2x + \varphi'(y) \]
    令其等于题目中给出的 $Q(x,y)$：
    \[ 2x + \varphi'(y) = 2x - 3y^2 \]
    在这个等式中，包含 $x$ 的项由于全微分的旋度判定条件完美相互抵消，化简得到只含有 $y$ 的微分关系：
    \[ \varphi'(y) = -3y^2 \]
    直接进行一元积分 $\varphi(y) = -y^3$（势函数积分常数可约定吸收到总常数内）。
    至此我们组合拼凑得到了隐式的通解势曲面簇 $u(x,y) = C$：
    \[ x^3 + 2xy - y^3 = C \]

    \textbf{方法二（配凑法：Method of Grouping）：}\\
    原方程不带条件直接强制拆分重组（按同类变量组合）：\\
    \[ (3x^2\d x - 3y^2\d y) + (2y \d x + 2x \d y) = 0 \]
    由于观察前几项极其显眼地满足标准全导数形式，我们马上可以将其包装回复合求导法则的形态：\\
    第一组直接就是普通的一元变量导数：$\d(x^3)$ 和 $-\d(y^3)$；\\
    第二组可以提出常数 $2(y \d x + x \d y) = 2\d(xy)$。\\
    得到完美的确切微分总和结构式：
    \[ \d(x^3) - \d(y^3) + 2\d(xy) = 0 \implies \d(x^3 - y^3 + 2xy) = 0 \]
    全微分等于零，立刻给出其原函数结构恒为一个不依赖时空演变的任意常数 $C$：
    \[ x^3 - y^3 + 2xy = C \]
    这远比复杂的两步“积分再求导对比”来得更加干净利落。
\end{solution}

\begin{example}
    利用配凑法求解：$x \d y - y \d x = (x^2+y^2) \d x$．
\end{example}
\begin{solution}
    观察到方程左端包含 $x\d y - y\d x$，右端有 $x^2+y^2$，这直接对应到 $\arctan(y/x)$ 的全微分核心结构。\\
    将等式两边同除以 $x^2+y^2$（此时设 $x^2+y^2\neq0$）：
    \[ \frac{x \d y - y \d x}{x^2+y^2} = \d x \]
    此时左边正是 $\d\left(\arctan\dfrac{y}{x}\right)$。\\
    于是方程化为 $\d\left(\arctan\dfrac{y}{x}\right) = \d x$。两边同时积分，得通解 $\arctan\dfrac{y}{x} = x + C$。
\end{solution}

\begin{exercise}
    求全微分方程$(x^3-y^2)\d x +(x^2y+xy)\d y=0$的通解．
\end{exercise}
\begin{solution}
    法一：\\
    设$u=y^2$，则原方程化为$\displaystyle \frac{\d u}{\d x}-\frac{2u}{x(1+x)}=-\frac{2x^2}{x+1}$，即
    \[\left(\frac{1+x}{x}\right)^2y^2+(1+x^2)=C\]
    法二：\\
    原方程化为
    \[
        \begin{array}{l@{\qquad}c}
            &\displaystyle x\d x+y \d y +y \frac{x\d y-y\d x}{x^2}=0\\[10pt]
            \Rightarrow & \displaystyle \dfrac{1}{2}\d(x^2+y^2)+y \d \left(\frac{y}{x}\right)=0\\[10pt]
            \Rightarrow & \displaystyle \dfrac{1}{2}\cdot \dfrac{\d(x^2+y^2)}{\sqrt{x^2+y^2}}+\dfrac{\left(\frac{y}{x}\right)}{\sqrt{1+\left(\frac{y}{x}\right)^2}}\d \left(\frac{y}{x}\right)=0\\[20pt]
            \Rightarrow & \displaystyle \d \sqrt{x^2+y^2}+\d \sqrt{1+\left(\frac{y}{x}\right)^2}=0\\[10pt]
            \Rightarrow & \displaystyle \sqrt{x^2+y^2}+\sqrt{1+\left(\frac{y}{x}\right)^2}=C\\
        \end{array}
    \]
    故通解为
    \[(1+x)\sqrt{x^2+y^2}=Cx\]
\end{solution}

\section{偏微分方程初步~(Introduction to PDEs)}
在常微分方程 (ODE) 中，未知函数只依赖于\textbf{一个}自变量（比如单纯随时间 $t$ 变化，或单纯随位置 $x$ 变化）。但在真实的物理世界中，一个基本量往往同时受时空影响，比如一根铁丝上的温度 $u(x,t)$ 既取决于你在铁丝上的哪个位置 ($x$)，又取决于现在是几点钟 ($t$)。这就引入了含有多个自变量偏导数的\textbf{偏微分方程 (PDE)}。

面对 PDE，初学者往往会感到无从下手。但请记住我们最核心的破局策略：\textbf{想尽一切办法把它退化成我们已经学过的一元 ODE。}

\subsection{三大基本方程（一般三维形式）}
对于物理世界中最一般的三维空间问题，刻画空间变化率（即各个方向局部凹凸程度的叠加）的绝佳数学工具是\textbf{拉普拉斯算子}（Laplacian），记为 $\nabla^2$ 或 $\Delta$：
\[ \nabla^2 u = \frac{\partial^2 u}{\partial x^2} + \frac{\partial^2 u}{\partial y^2} + \frac{\partial^2 u}{\partial z^2} \]
自然界最基础的 PDE 主要有以下三类：
\begin{enumerate}
    \item \textbf{波动方程~(Wave Equation):} $\dfrac{\partial^2 u}{\partial t^2} = a^2 \nabla^2 u$。描述声波、光波或引力波的传播。左边是时间二阶导（受力加速度），右边是空间二次偏导的汇总，这说明“一点处的场与其周围环境差值越大（越弯曲），它试图回弹的加速度就越大”。
    \item \textbf{热传导方程~(Heat Equation):} $\dfrac{\partial u}{\partial t} = \alpha^2 \nabla^2 u$。描述热量扩散或分子扩散。左边是时间一阶导（升降温速度）——说明“温度随时间的演化速度，正比于该点与其周围环境的温度差，总是向着熨平温差的方向进行”。
    \item \textbf{拉普拉斯方程~(Laplace Equation):} $\nabla^2 u = 0$。描述没有任何随时间变化 ($t$) 的最终平衡态，比如一个内部没有自由电荷的静电场势能分布，或者一块不均匀受热金属体达到恒温稳定后的空间三维温度分布。
\end{enumerate}

\subsection{降维打击：分离变量法 (Separation of Variables)}
既然我们目前只掌握了求解一元常微分方程 (ODE) 的本领，那求解偏微分方程最自然的物理直觉就是：假设时间演化和空间形态互不干涉。

我们大胆地猜测，二元函数 $u(x,t)$ 可以被拆解为两个一元函数的乘积：$u(x,t) = X(x) \cdot T(t)$。通过这种设定，我们能神奇地把原本复杂纠缠的 PDE，一刀劈成两个各自独立的 ODE。

\begin{example}[一维热传导方程初值问题]
    设长为 $L$ 的均质金属细杆绝热，两端温度坚壁冷却至 $0^\circ \text{C}$，试研究系统从任意初始分布的内部温度场 $u(x,t)$ 演化规律。
    \begin{align*}
        \textbf{热传导方程:} \quad & \frac{\partial u}{\partial t} = \alpha^2 \frac{\partial^2 u}{\partial x^2} \quad (0<x<L, \, t>0) \\
        \textbf{边界条件:} \quad & u(0,t) = 0, \quad u(L,t) = 0
    \end{align*}
\end{example}
\begin{solution}
    \textbf{第 1 步：剥离时空，转化为纯理论 O.D.E.}\\
    将 $u(x,t) = X(x)T(t)$ 代入主方程得 $X(x)T'(t) = \alpha^2 X''(x)T(t)$。\\
    重整化：$\dfrac{T'(t)}{\alpha^2 T(t)} = \dfrac{X''(x)}{X(x)}$。\\
    等式左边仅依赖时间 $t$，右边仅依赖空间结构 $x$，它们长久保持等于一个公用常数 $-\lambda$：
    \[ T'(t) + \alpha^2 \lambda T(t) = 0, \qquad X''(x) + \lambda X(x) = 0 \]
    
    \textbf{第 2 步：求解空间本征模态 (驻波动)}\\
    由杆两端结冰条件知 $X(0)=0$ 以及 $X(L)=0$。\\
    排除零解和指数发散解（$\lambda \le 0$），必定有特征值 $\lambda > 0$。令 $\lambda = \omega^2$，此时二阶常系数方程空间通解为谐振荡 $X(x) = C_1 \cos\omega x + C_2 \sin\omega x$。\\
    带入左端 $X(0) = C_1 = 0$；带入右端使得 $X(L) = C_2 \sin\omega L = 0$ 取非零，必要求 $\omega_n = \frac{n\pi}{L}$。\\
    故在空间上我们得到了\textbf{离散的本征基（正交函数族）}和\textbf{本征值}：
    \[ X_n(x) = \sin\left(\frac{n\pi x}{L}\right), \quad \lambda_n = \left(\frac{n\pi}{L}\right)^2, \; (n=1,2,3,\dots) \]
    
    \textbf{第 3 步：带回时间方程与完美级数收束}\\
    将得到的特有能级 $\lambda_n$ 送回时间常微分公式，得到解 $T_n(t) = e^{-\alpha^2 (n\pi/L)^2 t}$。\\
    由于原始偏微分方程全呈线性叠加律，最终温度场必然是全部离散本征场的线性级数组合：
    \[ u(x,t) = \sum_{n=1}^\infty b_n e^{-\alpha^2 \left(\frac{n\pi}{L}\right)^2 t} \sin\left(\frac{n\pi x}{L}\right) \]
    看那极速衰减的时间负指数项！随着时间 $t \to \infty$，所有内部高温波动终将迅速崩塌平息到绝对背景温度 $0$！\\
    系数常数 $b_n$ 怎么精确决定？我们惊喜地意识到，这只需要在初始时刻 ($t=0$) 算一个我们极为熟悉的——\textbf{傅里叶正弦级数展开}即可完美解决物理难题！
\end{solution}

\ifx\allfiles\undefined
\end{document}
\fi